\chapter{Decisions and repetition}
\MARSChapterHeading{Decisions and repetition}
\label{revision:decisions-and-repetition}

Every program so far has followed a fixed sequence from its first instruction to
its exit. A counter needs a choice: display another value, or stop. That choice
changes which instruction executes next. Repeating the same body with new register
values lets a short program process a whole range of values.

We will build a counter for MARS's right seven-segment display. Its data contains
sixteen segment masks, one for each hexadecimal digit from 0 through F. Two input
words select where to start and how many entries to display. The default inputs
select all sixteen. An empty range selects none, and must reach the exit without
reading a glyph or writing the display.

\section{Choose the next instruction}

The processor keeps the address of the next instruction in its \emph{program
  counter}, often abbreviated PC. During straight-line execution, the next address
follows the current instruction. A \emph{branch} can choose another instruction
address, specified in our source by a label.

Suppose \texttt{\$t1} holds the number of entries remaining. The instruction
\texttt{beqzc \$t1, done} branches to \texttt{done} if that register is zero.
Its mnemonic means ``branch if equal to zero, compact.'' If the value is nonzero,
execution continues with the following instruction. The assembler resolves the
label into the branch's relative destination; we do not calculate its encoded
distance by hand.

The branch reads \texttt{\$t1}; it does not decrement it. A counter needs separate
instructions to change its state. Similarly, the unconditional compact branch
\texttt{bc loop} always selects \texttt{loop}. The backward destination makes
repetition possible, but the register updates determine whether that repetition
will eventually end.

\section{Compact branches and the next source line}

These examples use MIPS64 Release 6 compact branches. A taken compact branch
transfers immediately: the instruction immediately after it is not executed on
that path. If \texttt{beqzc} finds zero, the next instruction executed is at
\texttt{done}. If it finds a nonzero value, the next instruction is the sequential
one after the branch. These two paths are called taken and not taken.

A conditional compact branch also has a \emph{forbidden slot}: when it is not
taken, the immediately following machine instruction cannot itself be a branch
or jump. Our program places \texttt{nop} there. This instruction performs no
useful data operation and does not transfer control, so either branch outcome
has a straightforward meaning. The instruction after the \texttt{nop} starts the
body only when entries remain.

A \texttt{nop} here is not a delay that runs before a taken branch. On the taken
path it is skipped. The unconditional \texttt{bc} has neither a delay slot nor a
forbidden slot. MARS also supports retained delayed branch instructions; copying
an older loop and changing only its branch mnemonic can change what executes.
Follow the semantics of the instruction you actually use.

\section{Build a two-way choice}

A branch can also choose between two computations without repeating either.
Suppose a flag is already in \texttt{\$t0}. We want \texttt{\$t1} to become
seven when the flag is zero, and nine otherwise. The following fragment belongs
inside a \texttt{main} routine after the flag has been prepared.

\MARSCodeLabel{Two paths with a shared continuation}
\begin{minted}{mars-idea}
    beqzc $t0, zero_case
    nop
    li $t1, 9
    bc joined
zero_case:
    li $t1, 7
joined:
    nop
\end{minted}

With a zero flag, the first branch selects \texttt{zero\_case}. It skips both
the safe \texttt{nop} and the assignment of nine. After assigning seven,
execution reaches \texttt{joined} by continuing sequentially. With a nonzero
flag, execution follows the \texttt{nop}, assigns nine and uses \texttt{bc} to
skip the other assignment. Both paths arrive at the shared continuation with
the chosen result in the same register.

The second \texttt{nop} marks a harmless instruction at the join for this small
experiment. In a complete program, useful work can continue there. The labels
organize the paths; they do not stop execution or automatically separate an
``else'' block. If we removed \texttt{bc joined}, the nonzero path would assign
nine and then continue into \texttt{zero\_case}, replacing it with seven.

Try flags zero, five and $-1$ in a copy of an earlier program's \texttt{main}
routine. Replace its calculation with this fragment and inspect \texttt{\$t1}
at \texttt{joined}, before exiting normally. Predict the result before each run.
The condition tests equality with zero, so both positive and negative nonzero
flags select nine. This pattern gives an assembly form of an ordinary two-way
conditional: a test, two paths, and a place where execution rejoins.

\section{A device address has an interface}

Chapter~\ref{revision:addresses-and-memory} used ordinary data memory for a bitmap.
Digital Lab Sim instead provides \emph{memory-mapped I/O}: selected addresses
represent device ports. A store to such a port supplies a value to a peripheral
model. The address, access width and written bits are part of its interface.

In the Default memory layout, the right seven-segment display port is at
\texttt{0x00000000FFFF0010}. The adjacent byte addresses serve other controls and
the left display. We must use a one-byte store to the right port. A word store
would touch more than that port, even if the source's low byte were correct.

The source uses \texttt{.eqv RIGHT\_DISPLAY} to give this constant a name.
\texttt{.eqv} defines an assembler substitution; it does not allocate a memory
object. \texttt{dli \$t7, RIGHT\_DISPLAY} loads the complete 64-bit constant into
a register. Like \texttt{li}, \texttt{dli} is an assembler convenience and can
expand into several instructions. Here it makes the intended address's zero
upper half explicit.

A byte written to this device is a \emph{segment mask}, with one bit per lighted
segment. It is not an ASCII character and is not the integer that the device will
somehow convert to a digit. With the standard segment assignment, the mask for
zero is \texttt{0x3F}, or binary \texttt{00111111}. Bits 0 through 5 are set;
the middle segment and decimal point are off.

\begin{center}
  \begin{tabular}{ll}
    \toprule
    Mask bit & Segment controlled    \\
    \midrule
    0        & Top horizontal bar    \\
    1        & Upper-right bar       \\
    2        & Lower-right bar       \\
    3        & Bottom horizontal bar \\
    4        & Lower-left bar        \\
    5        & Upper-left bar        \\
    6        & Middle horizontal bar \\
    7        & Decimal point         \\
    \bottomrule
  \end{tabular}
\end{center}

For digit one, only the two right bars are needed: bits 1 and 2 give
$2+4=6$, or \texttt{0x06}. A lookup table stores these designs as bytes, so the
loop does not need a different branch for every shape. Index zero selects mask
\texttt{0x3F}; index one selects \texttt{0x06}. The index and the selected mask
are different values with different jobs.

\section{Describe the loop's state}

The inputs satisfy $0\leq\texttt{first\_digit}\leq16$ and
$0\leq\texttt{digit\_count}\leq16-\texttt{first\_digit}$. These bounds keep every
requested entry inside the sixteen-byte table. They are preconditions of the
worked program, not checks it performs on arbitrary input.

We will retain the next table address in \texttt{\$t0}, the remaining count in
\texttt{\$t1}, and the next index in \texttt{\$t2}. A byte load obtains the mask
in \texttt{\$t3}, and \texttt{\$t7} retains the output port address. At the top
of the loop, the count decides whether the table address may be read.

After each store, advance the table address by one byte, advance the index by
one, and reduce the remaining count by one. Then return to the test. For the
initial range starting at zero with count sixteen, the first two iterations have
the following boundaries. The output history lists bytes in hexadecimal.

\begin{center}
  \begin{tabular}{llll}
    \toprule
    Boundary                          & Next index & Remaining & Output history    \\
    \midrule
    First test                        & 0          & 16        & Empty             \\
    After first store, before updates & 0          & 16        & \texttt{3F}       \\
    Second test                       & 1          & 15        & \texttt{3F}       \\
    Third test                        & 2          & 14        & \texttt{3F 06}    \\
    Final test                        & 16         & 0         & All sixteen masks \\
    \bottomrule
  \end{tabular}
\end{center}

The index immediately after a store still names the entry just displayed. At the
next test it names the next entry. State claims therefore need an execution
boundary, as they did for the earlier arithmetic and memory examples.

\section{The complete range counter}

The header connects Digital Lab Sim. The table and guard words are ordinary
static data; the display port is separate. The input words follow the table's
guard, so loading \texttt{first\_digit} and the next word obtains the two inputs.

\MARSCodeLabel{Display a bounded range of hexadecimal glyphs}
\begin{minted}{mars-idea}
# mars:tool digitalLab
.eqv RIGHT_DISPLAY 0x00000000ffff0010
.data
.align 2
glyphs:       .byte 0x3f, 0x06, 0x5b, 0x4f, 0x66, 0x6d, 0x7d, 0x07
              .byte 0x7f, 0x6f, 0x77, 0x7c, 0x39, 0x5e, 0x79, 0x71
glyph_guard:  .word 0x13579bdf
first_digit:  .word 0
digit_count:  .word 16
inputs_guard: .word 0x2468ace0

.text
.globl main
main:
    dli $t7, RIGHT_DISPLAY
    la $t0, glyphs
    la $t4, first_digit
    lw $t2, 0($t4)
    lw $t1, 4($t4)
    daddu $t0, $t0, $t2
loop:
    beqzc $t1, done
    nop
    lbu $t3, 0($t0)
port_write:
    sb $t3, 0($t7)
displayed:
    daddiu $t0, $t0, 1
    addiu $t2, $t2, 1
    addiu $t1, $t1, -1
iteration_done:
    bc loop
done:
    li $v0, 10
    syscall
\end{minted}

Save the program as \texttt{range-counter.asm} and select MIPS64 Release 6 with
the Default memory layout. It prints nothing. Open \emph{MARS Tools} and select
\emph{Digital Lab Sim} for the connected session. A full-speed Run may reach the
final F-shaped mask before you can see intermediate repaints; the program makes
no promise about animation speed.

The setup uses \texttt{daddu} to add the starting index to the table base as a
64-bit address calculation. Each table entry is one byte, so the index is also
the byte displacement. Inside the loop, \texttt{daddiu} advances that address by
one. The two \texttt{addiu} instructions update small 32-bit counts; their signed
word results fit throughout the admitted input range.

\section{Explain why it stops and stays in bounds}

Let $k$ be the number of completed iterations at a visit to \texttt{loop}.
At that point, the next index is $\texttt{first\_digit}+k$, the remaining count
is $\texttt{digit\_count}-k$, and exactly $k$ glyph bytes have been written in
order. This statement is a \emph{loop invariant}: it describes a relationship
that holds each time execution returns to the test.

Initially $k=0$, so no output has occurred. If the remaining count is nonzero,
the stated input bounds ensure the next index is less than sixteen. The byte
load is therefore inside the table. The body writes that byte once and updates
the registers to the state for $k+1$. The invariant holds again at the next test.

The remaining count starts nonnegative and decreases on each iteration. The
test occurs before the next load, so zero ends the loop without decrementing it
again. After the requested number of iterations, the next address may point one
byte past the selected range. That address is a useful boundary; it is not
permission to read another entry. In particular, starting at sixteen with count
zero computes a one-past-table address but never dereferences it.

If the count were decremented after zero, it would become negative rather than
signal a special ``finished'' condition by itself. The branch's location and the
updates together give the program its stopping behavior.

\section{Trace the choice in Debug}

Set a breakpoint at \texttt{loop}. With the default inputs, inspect the nonzero
remaining count and predict the not-taken path. Step over the branch and the
\texttt{nop}, then inspect the mask loaded into \texttt{\$t3}. At
\texttt{port\_write}, the store is still pending. After stepping to
\texttt{displayed}, the device has received \texttt{0x3F}, but the index and count
still have their old values.

At \texttt{iteration\_done}, inspect index one, remaining count fifteen and the
next table address. Continue to the second store and predict \texttt{0x06}.
For a shorter trace, change the inputs to start fourteen and count two. The
requested masks are \texttt{0x79} and \texttt{0x71}; the final index is sixteen.
Then try start sixteen and count zero. On its first visit, the test takes the
branch directly to \texttt{done}. The table body and the \texttt{nop} are skipped.

\section{Investigate a loop that looks almost correct}

A program that loads every table byte but stores only the last mask can leave the
same final display as the complete counter. It has not performed the requested
sequence. A device interaction is an ordered history of transactions, not just
a final register or picture. This is why the project's checks observe the count,
order, width, address and values of actual stores.

If the output skips every second entry, inspect the table-address increment.
An increment of two advances two bytes. If the output never stops, inspect whether
the branch reads the same count that the body decrements. If an empty range writes
one mask, check whether a load/store was placed before the first zero test.
Each symptom suggests a specific state or execution boundary to inspect.

Putting another branch immediately after a not-taken \texttt{beqzc} is a different
fault: it violates the forbidden-slot rule. A label between those instructions
does not create an intervening instruction. Keep the safe \texttt{nop} while
learning to follow both paths.

\section{Write a different loop with the same contract}

In a separate version, retain an exclusive end address rather than a remaining
count. Stop when the current address equals that end. Use a compact equality
comparison such as \texttt{beqc}, and keep its following instruction safe. Derive
the end address from the supplied first index and count. The empty-range case
must still perform no table load or port store.

Test an empty range, one entry, the full range and a range ending exactly at
sixteen. Preserve the table and both guards. Your registers and labels may differ
from the worked program; the requested output transactions must agree.

For another experiment, alter the glyph design for one in the table. Predict the
segments from its bits, then display just that entry. This changes the table's
meaning without changing the loop's control structure. The elective MarsBot lab
uses the same ideas of state, repetition and device writes to control movement.
